\subsection{Descomposición geométrica bajo inversión de hoja}
\label{sec:ii09-descomposicion-paridad}

En la coordenatización angular, la inversión de hoja conserva \(a,d\) e
invierte \(u\). El vector expresado admite la descomposición
\[
 \boldsymbol\theta_+=a c_A+u c_h+d c_\Delta,
 \qquad
 \boldsymbol\theta_-=a c_A-u c_h+d c_\Delta.
\]
\begin{proposition}[Componentes par e impar del transporte angular]
\label{prop:ii09-paridad-angular}
Los dos vectores satisfacen
\[
 \frac{\boldsymbol\theta_++\boldsymbol\theta_-}{2}
       =a c_A+d c_\Delta,\qquad
 \frac{\boldsymbol\theta_+-\boldsymbol\theta_-}{2}=u c_h.
\]
La contribución que cambia de signo ocupa la recta
\(\operatorname{span}(c_h)\); la componente conservada ocupa
\(\operatorname{span}(c_A,c_\Delta)\).
\end{proposition}
\begin{proof}
La suma cancela los términos \(u c_h\), y la diferencia cancela
los otros dos. Las columnas son linealmente independientes, pues
\(\det\mathsf M_{\rm CKM}=-3857/6\). Por ello las dos
componentes forman una descomposición directa.
\end{proof}

La misma inversión actúa sobre la evaluación de Barbero--Immirzi:
intercambia \(q_+\) y \(q_-\), e invierte la componente bilineal
en \(A\) y \(C^*\). En consecuencia, el funcional de incidencia
es impar. Un observable simétrico en ambos canales conserva su
valor. Esta diferencia de paridad expresa cómo una misma
operación sobre la estructura produce respuestas distintas según
el lector: orientación areal, información simétrica y vector de
mezcla conservan información diferente del mismo antecedente.
